% Authoring template, not a completed coursework report.
% Replace every \todo{...} with your own verified work; remove instructional prose.
% Keep the complete report, references and disclosures within five pages.
% The main document supplies the title, member details, bibliography and disclosures.

\section{Problem and motivation}
\todo{State the plant disease recognition problem, its practical motivation, and the
scope of the nine-host coursework subset. Identify the selected option and explain
what a successful result would establish. Avoid claims unsupported by your results.}

\ifnum\ReportOption=2
\todo{State a falsifiable hypothesis before training: explain how binary lesion
supervision could improve rejection of unfamiliar host--disease combinations while
preserving known-class recognition. Define and justify an acceptable reduction in
known-class macro-F1, and specify what evidence would support, contradict, or leave
the hypothesis unresolved.}
\else
\todo{Define the closed-set prediction task: a single RGB image is mapped to a
host--disease class. Explain why visually similar symptoms, small lesions, background
variation and imbalance make this task challenging.}
\fi

\section{Related work}
The required dataset release is PlantSeg v7 \citep{wei2025plantsegv7}; the associated
data descriptor provides the dataset context \citep{wei2026plantseg}.
\todo{Add the relevant methods you actually read to references.bib. Explain how their
data, supervision, inference inputs and evaluation protocols relate to this study.
Check every bibliographic record and cite sources for substantive borrowed ideas.}

\ifnum\ReportOption=2
\todo{Discuss at least three closely related methods, including unknown plant-disease
recognition and lesion-guided or grounded recognition. Record the date and scope of
the current literature search. Distinguish rejecting unknown classes from naming
unseen classes or abstaining on healthy images. Explain whether this study is a
replication, extension or new comparison; substantiate any novelty claim.}
\fi

\section{Data and protocol}
\todo{Use only PlantSeg v7 (25 November 2025), retaining apple, banana, citrus, corn,
potato, soybean, tomato, wheat and zucchini. Derive actual class identifiers and
counts from the downloaded metadata. Treat each host--disease pair as one class;
preserve labels and the supplied approximately 70/10/20 train/validation/test split.}

\todo{Reference the frozen manifest and report per-partition class counts, unreadable
files, duplicate checks and documented exclusions. Explain how identical images and
possible near-duplicate leakage were handled. Fit dataset-specific preprocessing on
training data only; keep augmentations with their source training images. State the
fixed class order, treatment of undefined metrics, and handling of unsupported classes.}

\ifnum\ReportOption=2
\todo{Document both frozen class-role mappings. Citrus classes are development
unknowns and potato classes are test unknowns. For each other host, Assignment A
reserves the first/second alphabetically ordered disease labels as development/test
unknowns; Assignment B reserves the second/third. Train on known classes only and
discard unused reserved-class images. Report familiar-host unknowns separately from
the unfamiliar-host control, including shared versus unseen disease labels where
sample sizes permit.}
\fi

\section{Method}
\todo{Describe the architecture, initialisation or exact pretrained-weight identifier,
RGB preprocessing, augmentation, loss, optimiser, learning rate, training budget and
stopping rule. Justify resizing/cropping. Record seeds, dependency versions, hardware,
training time and checkpoint selection. Use deterministic validation/test transforms,
evaluation mode and disabled gradients. Do not use PlantSeg-trained or overlapping
PlantWild-trained checkpoints, class-bearing filenames, or metadata as predictive inputs.}

\ifnum\ReportOption=2
\todo{Explain the lesion-supervision mechanism and why it might succeed or fail.
Use binary masks from known training classes only; all evaluated predictions must
use RGB only. Define a matched classifier without lesion supervision and control
backbone, initialisation, preprocessing, augmentation and training/search budgets.
Remove each claimed additional component in an ablation.}
\else
\todo{Describe a majority-class predictor and a basic CNN. Define at least two
controlled improvements, changing one factor at a time where possible. State how
validation macro-F1 selects the final configuration and how all choices are frozen
before final testing.}
\fi

\section{Experiments and results}
\todo{Replace all dashes in Table~\ref{tab:results} with verified notebook-generated
results, with sample sizes and uncertainty where required. A dash is a placeholder,
not zero or an omitted experiment. Report macro-F1 (宏平均 F1) consistently.}

\begin{table}[htbp]
\centering
\small
\ifnum\ReportOption=2
\caption{Results template for one class assignment. All entries are placeholders;
no experiments or measured results are represented.}
\label{tab:results}
\begin{tabularx}{\linewidth}{@{}lXXX@{}}
\toprule
Model & Familiar-host AUOSCR & Known-class macro-F1 & Unknown AUROC \\
\midrule
Matched baseline & -- & -- & -- \\
Lesion supervision & -- & -- & -- \\
Component ablation & -- & -- & -- \\
\bottomrule
\end{tabularx}
\else
\caption{Test-results template. All entries are placeholders; no experiments or
measured results are represented.}
\label{tab:results}
\begin{tabularx}{\linewidth}{@{}lXXX@{}}
\toprule
System & Macro-F1 & Top-1 accuracy & Balanced accuracy \\
\midrule
Majority predictor & -- & -- & -- \\
Basic CNN & -- & -- & -- \\
Controlled change 1 & -- & -- & -- \\
Controlled change 2 & -- & -- & -- \\
\bottomrule
\end{tabularx}
\fi
\end{table}

\ifnum\ReportOption=2
\todo{Run three training seeds for each of two class assignments: six runs per model
configuration, with independent training for each assignment. Report seed means,
standard deviations and paired differences separately for each assignment. Compare
maximum-softmax, maximum-logit and energy scores; choose the primary rule on baseline
validation data and apply the same rule to both models. Report familiar-host AUOSCR,
known-class macro-F1/accuracy, unknown AUROC and separate unfamiliar-host results.}
\todo{Select checkpoints by known-class validation macro-F1 and configurations by
familiar-host validation AUOSCR subject to the predefined tolerance. Freeze both
assignments' choices before either test evaluation. Set the operating threshold using
known validation images only at approximately 95\% acceptance; report test known
acceptance, correct-and-accepted rate and unknown false acceptance.}
\else
\todo{Report final test metrics for the baseline and improved systems, per-class
precision/recall/F1/support, and a confusion matrix. Retain the complete matrix in the
notebook if the report shows a readable subset. Explain the effects of the controlled
changes using the actual evidence.}
\fi

\begin{figure}[htbp]
\centering
\fbox{\parbox[c][23mm][c]{0.9\linewidth}{\centering
Figure placeholder / 图示占位\\
No data and no measured results are shown.\\
Replace with an informative notebook-generated visual.}}
\caption{Visual-evidence placeholder, not a measured result. Replace both this box
and its caption with a numbered, labelled figure generated by the submitted notebook.}
\label{fig:diagnostic}
\end{figure}

\todo{Replace Figure~\ref{fig:diagnostic} with a useful result: for example, a confusion
heatmap, checkpoint-marked learning curves, labelled prediction failures, or Option 2
OSCR curves. Give axis/colour labels, clear class labels and relevant sample counts.
The submitted notebook must regenerate every data-derived figure and results table.}

\section{Interpretation and limitations}
\todo{Explain representative successes and failures, class imbalance, possible leakage,
data/domain limits and computational constraints. Separate observed results from
speculation and avoid broad claims unsupported by the evaluated groups. Summarise
what the evidence establishes and one justified next investigation.}
\ifnum\ReportOption=2
\todo{Evaluate the pre-stated hypothesis against paired AUOSCR and known-class
macro-F1 differences, variation across seeds/assignments and unknown groups. Mixed or
negative evidence is valid. A potato-only unfamiliar-host control does not justify
claims about general performance on unfamiliar plant hosts.}
\fi

% References and AI/team disclosures are rendered centrally by the main template.
